% New research after the restored quadratic count; separate from frozen assembly.
\begin{proposition}[Proved under stated hypotheses: projective representations with bounded central obstruction]
\label{D:bounded-centre-projective-growth}
Let $Q$ be a profinite group and fix a prime $p$. Assume that
$H^2_{\mathrm{cont}}(Q,\mathbf Z/d\mathbf Z)$ is finite for every
positive integer $d$ prime to $p$, with trivial action. Assume also
that every open normal subgroup $H\subset Q$ has a constant $c_H$
such that
\[
 \#\bigl(\operatorname{Hom}_{\mathrm{cont}}
 (H,\operatorname{GL}_r(\mathbf F_{p^f}))/\mathrm{conj}\bigr)
 \le p^{c_Hfr^2}\qquad(r,f\ge1).
\]
For every fixed integer $B\ge1$ there is a constant $C_B$ with the
following property. Put $q=p^f$ and $d=\gcd(r,q-1)$. If $d\le B$,
then
\[
 \#\bigl(\operatorname{Hom}_{\mathrm{cont}}
 (Q,\operatorname{PSL}_r(\mathbf F_q))/
       \operatorname{PGL}_r(\mathbf F_q)\bigr)
 \le q^{C_Br^2}.
\]
Conjugation in the displayed quotient is through the natural normal
subgroup $\operatorname{PSL}_r(\mathbf F_q)\subset
\operatorname{PGL}_r(\mathbf F_q)$. In particular the number of
distinct kernels of epimorphisms to this simple group is bounded by
the same expression whenever the group is nonabelian simple.
\end{proposition}
\begin{proof}
We first construct one subgroup independently of $r$ and $f$. A
continuous normalized two-cocycle from $Q\times Q$ to a finite
discrete group is constant on the pairs of cosets of some open
normal subgroup $N$. Indeed compactness reduces its clopen fibres
to finitely many rectangles, and a common open normal subgroup
refines all their factors. Thus the cocycle factors through
$(Q/N)^2$. Its restriction to $N\times N$ is its value at $(1,1)$,
which is zero. Every continuous degree-two class with finite
coefficients is therefore killed by restriction to an open normal
subgroup; this is also the continuous-cochain calculation in
\citet[II, Proposition 4.2, p.~87]{MilneCFT2020D}.

There are only finitely many prime-to-$p$ integers $d\le B$, and by
hypothesis only finitely many classes for each. Intersect the open
normal subgroups just obtained. We obtain a single open normal
subgroup $H_B\subset Q$ such that
\[
 H^2(Q,\mathbf Z/d\mathbf Z)\longrightarrow
 H^2(H_B,\mathbf Z/d\mathbf Z)
 \quad\hbox{is zero for every such }d.
\]
We require neither an effective bound for its index nor one uniform
in $B$.

The central extension
\[
 1\longrightarrow\mu_d(\mathbf F_q)
 \longrightarrow\operatorname{SL}_r(\mathbf F_q)
 \longrightarrow\operatorname{PSL}_r(\mathbf F_q)
 \longrightarrow1
\]
has cyclic kernel of order $d$, which is prime to $p$. The pullback
of this extension by any continuous representation $\rho$ of $Q$
has a class in $H^2(Q,\mathbf Z/d\mathbf Z)$, after choosing an
identification of its centre with that cyclic group. The class
restricts to zero on $H_B$. A continuous cochain trivializing the
restricted cocycle supplies a continuous homomorphism
\[
 \widetilde\rho_H:H_B\longrightarrow
                   \operatorname{SL}_r(\mathbf F_q)
\]
lifting $\rho|_{H_B}$. This is a lift over the original finite
field; no scalar-field extension is involved.

Every conjugacy class of a restricted projective representation
therefore occurs as the projectivization of a linear representation
of $H_B$. Hence the number of possible restricted classes is at
most $q^{c_{H_B}r^2}$. Put $t=[Q:H_B]$ and choose $t$ coset
representatives, including $1$. After choosing one representative
of a restricted projective conjugacy class, any extension to $Q$
is determined by its values on these coset representatives. There
are at most $|\operatorname{PGL}_r(\mathbf F_q)|^{t-1}$ possible
tuples. Relations can only reduce their number. Conjugation needed
to fix the restriction causes no problem, since every element of
$\operatorname{PGL}_r(\mathbf F_q)$ is represented by a matrix in
$\operatorname{GL}_r(\mathbf F_q)$.

Since $|\operatorname{PGL}_r(\mathbf F_q)|\le q^{r^2}$, the total
is at most $q^{(c_{H_B}+t-1)r^2}$. Finally conjugate homomorphisms
have the same kernel, so the bound for kernels follows.
\end{proof}

\begin{corollary}[Proved: bounded-centre type-A simple quotients]
\label{D:surface-bounded-centre-simple-product}
Let $Q=Q(S,T)$ be the original-boundary tame group of a normal proper
surface over an algebraically closed field of characteristic $p>0$.
Fix $B$. The quotient obtained by taking the product
of all distinct nonabelian simple quotients
$\operatorname{PSL}_r(\mathbf F_{p^f})$ with
$\gcd(r,p^f-1)\le B$ is topologically finitely generated. The ranks
$r$ are not bounded in this assertion.
\end{corollary}
\begin{proof}
Theorem~\ref{D:surface-tame-cross-characteristic-H2} implies the
finiteness assumption of
Proposition~\ref{D:bounded-centre-projective-growth}. An open normal
subgroup $H\subset Q$ corresponds to a finite etale cover $U_H\to U$
of $U=S\setminus T$. Normalize $S$ in this cover to obtain a normal
proper surface $S_H$, and put $T_H=S_H\setminus U_H$. At each
original boundary divisor the cover is tame. The intersection
property for wild inertia shows that every finite representation
of $H$ is tame at the original divisors of $T_H$: its monodromy is
inflated from $H$, and the new wild inertia maps to the old killed
wild inertia. Equivalently there is a surjection
$Q(S_H,T_H)\twoheadrightarrow H$, which is all we need.
Theorem~\ref{D:tame-global-quadratic-representation-bound} therefore
supplies the hereditary quadratic bound used in the proposition.

For these simple groups the order formula
\[
 |\operatorname{PSL}_r(\mathbf F_q)|
 =\frac{q^{r^2-1}}{\gcd(r,q-1)}
       \prod_{i=2}^r(1-q^{-i})
\]
gives $|\operatorname{PSL}_r(\mathbf F_q)|\ge q^{r^2/4}$ for
$r\ge2$. Indeed the product is at least
$1-\sum_{i\ge2}2^{-i}=1/2$, while $2r\le2^r\le q^r$ and
$r^2-r-1\ge r^2/4$ for $r\ge2$. Thus Proposition~\ref{D:bounded-centre-projective-growth}
bounds the multiplicity of every such simple quotient by a fixed
power of its order.

The standard generating-tuple estimate for a nonabelian simple
group, used already in Proposition~\ref{D:bounded-matrix-simple-product},
then gives one number of generators for all the corresponding finite
products. This is the estimate in the proof of Proposition~2.6(ii),
printed pp.~12--13, of \citet{LucchiniMarionTracey2020D}.
Distinct simple quotient kernels give independent factors by
Goursat's lemma, so the resulting profinite quotient is their full
product. Compactness of the finite generating-tuple conditions
gives topological finite generation of that product.
\end{proof}

\begin{remark}[Scope of the remaining central gate]
The fixed subgroup $H_B$ depends on $B$. This argument supplies no
uniform bound as $\gcd(r,q-1)$ grows. In particular it does not
prove an all-rank projective representation bound with one constant
for every type-A simple group. A uniform count for twisted unit
bundles, or another controlled central-lifting argument, is still
needed for that stronger conclusion. Nor does a bound on simple
direct-product quotients by itself handle general nonabelian crowns.
\end{remark}
